\documentclass[11pt]{article}
\usepackage[a4paper,margin=25mm]{geometry}
\usepackage{amsmath,amssymb,booktabs,longtable,hyperref}
\newcommand{\Lc}{\textbf{(c)}}
\newcommand{\Ld}{\textbf{(d)}}
\emergencystretch=2em
\title{Independent audit of the Boussinesq fluctuation determinant foundations}
\author{Project audit}
\date{4 November 2026}
\begin{document}
\maketitle

\paragraph{Scope and result.}
This note audits the logical passage from the exact phonon solutions and their
phase integrals to a physical one-loop two-soliton or soliton--defect amplitude.
The good-equation and potential-field conventions of the repository are retained.
No regularisation, asymptotic state convention or exchange phase is selected.
The nonconstant continuum integral has substantial positive evidence. Its
identification with the full determinant, including zero-mode measure and an
absolute phase, remains incomplete \Ld. The label \Ld\ here denotes an
unestablished physical conclusion, not a counterexample to the continuum
integral. All new numerical and symbolic identities below are checked by
\texttt{scripts/audit\_\allowbreak determinant\_\allowbreak foundations.py}, at 65-digit working precision;
raw results are saved in \texttt{scripts/audit\_determinant\_results.json}.

\section{What the cited determinant argument supplies}

Faddeev--Korepin (FK) \cite{FK} derives the relevant one-loop determinant ratio
for relativistic scalar fields, with the Klein--Gordon fluctuation operator.
In \S5.4, printed p.~62, the constant factor $Z$ is explicitly left to
supplementary requirements, with crossing symmetry as an example. Equation
(4.28), printed p.~66, still contains $Z$. This is a statement of the source's
scope \textbf{(a)}, not a Boussinesq theorem. The source's Appendix~6 derives its
spectral density from a second-order Schr\"odinger equation, and Appendix~5,
p.~81, uses localization of discrete modes to separate their Wronskians.

The Boussinesq operator has fourth-order spatial derivatives and derivative
coupling to the background. A canonical second-order time derivative ensures
the form of the temporal symplectic product; it does not alone establish
completeness, the resolvent, the infrared domain, or the source's entire
spectral-density derivation. Thus ``apply verbatim'' in
\texttt{oneloop\_phase.tex} is too strong. The three-soliton tau linearisation
provides an exact continuum family with diagonal transmission. Conditional on
its completeness in the relevant scattering space, this supports a vanishing
continuum mixing kernel. It does not itself construct the complete projected
resolvent with collective-coordinate directions removed \Ld.

\paragraph{A co-moving sign correction.}
The displayed co-moving ODE in Phase~2e has a typographical sign error.
Direct symbolic differentiation gives \Lc
\[
 \partial_t^2\{e^{-i\widetilde\omega t}f(x-vt)\}
 =e^{-i\widetilde\omega t}
 (-\widetilde\omega^2f+2iv\widetilde\omega f'+v^2f'').
\]
Its first-derivative term therefore has a plus sign. The subsequent quartic
in that report uses the correct sign; the lab-frame PDE checks are unaffected.

\paragraph{Positive control: the potential-field density.}
A direct calculation does support the one-soliton density step \Lc. Write
$k=\sin\alpha$, $v=\epsilon\cos\alpha$, $\omega=|q|\sqrt{1+q^2}$,
$h=(1+\tanh(k\xi/2))/2$, $A=e^{i\theta(q)}$ and
\[
 g(\xi)=1+(A-1)\left(h+\frac{h'}{iq}\right),\qquad
 \psi(x,t)=\omega^{-1/2}g(x-vt)e^{iqx-i\omega t}.
\]
This is the potential-field tau mode, normalized to temporal plane-wave
Wronskian density $-2i$. Direct spatial quadrature gives
\begin{equation}\label{eq:density}
 \frac1\omega\int_{\mathbb R}\!
 \left[-2i\omega(|g|^2-1)-v(\bar g g'-g\bar g')\right]d\xi
 =-2i\,\theta'(q). \tag{c}
\end{equation}
Four checks include both soliton orientations and both phonon directions.
This checks the normalization that is implicit in the quoted $-4\pi i$
continuum Wronskian. It is not a check of the full time-dependent determinant.

\section{Endpoint constants and exchange}
For the phase continuous on each half-line and zero at both infinities,
the light-branch endpoints are \Lc
\begin{equation}
 \theta(0^\pm;\alpha,\epsilon)=\pi(\pm1+\epsilon),
 \qquad 0<\alpha<\pi/3. \tag{c}
\end{equation}
Define the ordinary half-line integrals (without a delta function at zero)
\[
 I_{12}=\frac1{4\pi}\int_{\mathbb R\setminus\{0\}}\theta_1\theta'_2\,dq,
 \qquad J_{12}=\frac{I_{12}-I_{21}}2.
\]
Their integration-by-parts identity, for all four direction combinations, is
\begin{align}
 I_{12}+I_{21}&=-\frac\pi2(\epsilon_1+\epsilon_2),\tag{c}\\
 J_{12}-I_{12}&=\frac\pi4(\epsilon_1+\epsilon_2),\tag{c}\\
 I(\alpha,\epsilon;\alpha,\epsilon)&=-\epsilon\frac\pi2.\tag{c}
\end{align}
The \texttt{oneloop\_phase.tex} same-direction expression $-\pi$ is correct for
$++$ and has the wrong sign for $--$. In particular, the later claims that raw
$I$ is antisymmetric and vanishes on the diagonal are false. The displayed
identities are checked by direct quadrature of a phase reconstructed from the
KP transmission ratio, independently of the closed dilogarithm formula.

The endpoint constants are independent of the amplitudes inside each regular
chamber. Consequently the FK gradient cannot distinguish $I$ from $J$ or fix a
further chamber constant. Phase~2e itself acknowledges the $Z$ ambiguity and
records the PI's decision to display both functions. A later preference for
one function because it closes a KP identity is a normalization criterion,
not a calculation of that determinant constant. In particular, a phase
integral evaluated after both particle labels are swapped is not automatically
the inverse physical collision: the ordering of incoming trajectories, the
phase of the asymptotic state basis, and any exchange convention must also be
specified. Neither the raw exchange failure nor algebraic closure alone
settles physical unitarity \Ld.

For the defect integral, setting the defect phase to zero at $q=0^\pm$ makes
the integration-by-parts term vanish. That is a precise statement about a
chosen branch. The absence of a translation mode does not select a logarithm
branch: $z_{\rm ph}(0)=1$ fixes it only modulo $2\pi$. No new defect branch is
chosen in this audit. A physical fusion rule additionally requires normalized
intermediate defect states and fusion amplitudes; continuum chord identities
alone do not determine those data \Ld.

\section{Collective coordinates and the Maslov assertion}
At $t=0$, choose the repository potential gauge
\[
 \phi(x;\alpha,X)=k\{1+\tanh[k(x-X)/2]\},\qquad \pi=-v\phi_x.
\]
The single-soliton symplectic pairing verifies the canonical momentum \Lc:
\begin{equation}
 \Omega(\partial_\alpha,\partial_X)
 =\int(\pi_\alpha\phi_X-\pi_X\phi_\alpha)dx
 =\frac{dP}{d\alpha},\qquad
 P=\frac23\epsilon\sin^3\alpha\cos\alpha.\tag{c}
\end{equation}
However the amplitude variation is not a spatially localized potential mode:
\begin{equation}
 \phi_\alpha(-\infty)=0,\qquad
 \phi_\alpha(+\infty)=2\cos\alpha.\tag{c}
\end{equation}
The variation is a generalized collective-coordinate direction, including a
secular component at nonzero time, rather than a second ordinary oscillatory
bound state. Statements about localized $\delta u$ cannot be transferred to
the potential-field determinant without tracking this tail.

For two widely separated asymptotic profiles, soliton~1 on the left and
soliton~2 on the right, retaining these tails gives \Lc
\begin{equation}\label{eq:mixed}
 \Omega(\partial_{\alpha_1},\partial_{\alpha_2})
 \longrightarrow -M'_1\Pi'_2,\qquad
 M=2\sin\alpha,\quad \Pi=-2\epsilon\sin\alpha\cos\alpha.\tag{c}
\end{equation}
The check evaluates the spatial integral at separations 200 and 400.
This is an asymptotic-profile calculation, not a substitute for the exact
collision modes. It establishes that separation alone does not make the
mixed block vanish. Darboux coordinates may remove these cross terms, but
the corresponding zero-mode basis, normalization and measure must actually
be supplied. Conservation of a Wronskian establishes time independence;
it does not by itself establish equality to a product of one-body factors.

\paragraph{What the large-time Hessian does show.}
The checked derivative is \Lc
\begin{equation}
 \frac{dv}{dP}=-\frac{3}{2\sin3\alpha}<0,
 \qquad 0<\alpha<\pi/3.\tag{c}
\end{equation}
For fixed parameters away from degeneracies and resonance lines, a real
symmetric finite correction to $T\,\mathrm{diag}(dv_j/dP_j)$ cannot change its
negative signature at sufficiently large $T$. A matrix control in the script
keeps the off-diagonal terms; the diagonal entries quoted as ``eigenvalues''
in Phase~2e are not generally eigenvalues. The valid signature observation
concerns the endpoint stationary-phase wave-packet factor. It is not a
calculation of spectral flow of the field fluctuation operator or of the
branch of its square root. It therefore does not establish a full relative
Maslov index or determine $Z$ \Ld. The cited Phase~2e script contains no
Wronskian, measure, Hessian or Maslov test despite being cited for that claim.

A triangular coordinate shift with unit finite-dimensional Jacobian and a
classical function of commuting momenta do rule out a particular elementary
ordering problem. They are not an evaluation of the projected functional
measure or its global phase. Thus the claimed absence of every additional
one-loop topological/measure contribution is unproved \Ld.

\section{Localized defect modes and the physical domain}
Phase~2g asserts that localized defect modes cancel because the defect has
identical endpoint frequencies before and after a soliton crosses. Neither
of its cited scripts evolves the localized modes through the collision.
Endpoint frequencies alone do not prove cancellation. As a logical control,
an oscillator with frequency $\omega$ before and after a finite interval of
length $T$ and frequency $W$ inside has the exact negative-frequency
coefficient \Lc
\begin{equation}
 \beta=i\frac{W^2-\omega^2}{2\omega W}\sin(WT).\tag{c}
\end{equation}
At $(\omega,W,T)=(1,2,1)$ it is nonzero. The script derives this coefficient
from the exact transfer matrix and verifies the symplectic normalization.
This counterexample refutes the general implication used in the text;
it does \emph{not} assert that this mixing occurs in the integrable defect.
That would require its actual localized scattering matrix. In the strictly
stable regime with no localized defect modes this particular obstruction
is absent; at thresholds and unstable vacua it cannot be dismissed by the
same endpoint argument \Ld.

The most defensible domain of a putative elastic two-soliton interpretation
is the light branch with a nonsingular real two-soliton background, distinct
velocities, and parameters away from cusp/resonance thresholds. Continuum
integrals can remain finite beyond this domain without defining asymptotic
stable scattering states. The tau construction supplies candidate discrete
modes and the continuum phase winding, but a complete spectral theorem for
the fourth-order operator pencil (with its zero-frequency threshold and
potential-field domain) is not demonstrated by the existing finite samples.
A ``Levinson-type count'' should remain distinguished from such a theorem
\Ld. The defect stability formula itself is audited separately.

\section{Claims and checks}
\begin{longtable}{p{0.47\linewidth}p{0.08\linewidth}p{0.36\linewidth}}
\toprule Claim & Label & Evidence\\\midrule
Normalized potential-mode density identity & (c) & Four 65-digit spatial quadratures\\
Orientation-dependent endpoint and exchange constants & (c) & Four channel quadratures, diagonal checks\\
Single and mixed asymptotic symplectic pairings & (c) & Direct field integrals and symbolic momentum derivative\\
Endpoint Hessian signature at large time & (c), limited & Full matrix control; no field spectral-flow claim\\
Equal localized frequencies need not imply cancellation & (c) & Exact oscillator transfer matrix\\
Full Boussinesq determinant and collective measure & (d) & Missing complete projected resolvent and normalization\\
Absolute one-loop phase, exchange/fusion constants & (d) & FK source retains $Z$; no independent anchor\\
\bottomrule
\end{longtable}

\section{Failures recorded}
\begin{enumerate}
\item The co-moving ODE has the wrong first-derivative sign; its subsequent
quartic and lab-frame PDE checks have the correct sign.
\item Prior orientation-independent same-direction boundary formula fails for
$--$: its exchange sum is $+\pi$.
\item Prior raw antisymmetry and zero diagonal claims fail; the diagonal is
$-\epsilon\pi/2$.
\item Localization-only cancellation of amplitude-mode Wronskians fails for
the potential-field variations: their mixed asymptotic pairing is nonzero.
This does not disprove cancellation after a correctly constructed Darboux
measure; that calculation is missing.
\item Equal asymptotic localized frequencies do not imply determinant
cancellation, as the exact pulse control demonstrates.
\item The first Hessian illustration incorrectly required an eigenvalue to
differ from a diagonal entry by more than 1; this arbitrary bound failed. It
was replaced by the exact symbolic characteristic-polynomial identity at each
diagonal entry, which equals minus the square of the off-diagonal entry and
is therefore nonzero in the test. No tolerance was weakened. The checks do not evaluate a functional determinant or prove
completeness and are not reported as doing so.
\end{enumerate}

\begin{thebibliography}{9}
\bibitem{FK} L.~D.~Faddeev and V.~E.~Korepin,
\emph{Quantum theory of solitons}, Physics Reports \textbf{42} (1978), 1--87,
\href{https://doi.org/10.1016/0370-1573(78)90058-3}{DOI:10.1016/0370-1573(78)90058-3}.
The publisher's entry and DOI were verified. Relevant full-text passages in
\S5.4 and Appendices~3--6 were inspected from the local PDF
\texttt{Faddeev and Korepin - 1978 - Quantum theory of solitons.pdf};
this audit does not claim to have read the entire article.
\end{thebibliography}
\end{document}
